\section{Determinant}
\label{sec:determinant}

\noindent\textbf{Subject:} Linear Algebra for ML (3Blue1Brown) \\
\textbf{Status:} growing

\subsection*{Summary}
The determinant measures \textbf{how much a linear transformation scales areas}. It can be negative (flipping orientation) and zero means the transformation squishes space into a lower dimension.

\subsection{Core Idea}
\begin{itemize}[nosep]
  \item \emph{How much are areas scaled?}
  \item The linear transformation scales its area
  \item The special scaling factor by which a linear transformation changes any area is called the \textbf{determinant}
\end{itemize}

\subsection{Formula}

\paragraph{$2\times2$ Matrix:}
\[
\det\left(\begin{bmatrix} a & b \\ c & d \end{bmatrix}\right) = ad - bc
\]

\paragraph{$3\times3$ Matrix:}
\[
\det\left(\begin{bmatrix} a & b & c \\ d & e & f \\ g & h & i \end{bmatrix}\right)
= a \det\begin{bmatrix} e & f \\ h & i \end{bmatrix}
- b \det\begin{bmatrix} d & f \\ g & i \end{bmatrix}
+ c \det\begin{bmatrix} d & e \\ g & h \end{bmatrix}
\]

\subsection{Notes}
\begin{notebox}
\begin{itemize}[nosep]
  \item The determinant \textbf{can be negative} because of the scaling (flipping orientation)
  \item \textbf{Question:} $\det(M_1 M_2) = \det(M_1) \cdot \det(M_2)$
  \begin{itemize}[nosep]
    \item $\rightarrow$ Applying two transformations in sequence multiplies their area-scaling factors
    \item $\rightarrow$ The determinant is that scaling factor
  \end{itemize}
\end{itemize}
\end{notebox}

\subsection*{Related Topics}
\begin{itemize}[nosep]
  \item Section~\ref{sec:matrices} --- Matrices
  \item Section~\ref{sec:linear-systems} --- Linear System of Equations
\end{itemize}

\bigskip
\noindent\textit{Tags: linear-algebra, determinant, math, phase-1}
